\section{Los números enteros y la sustracción}

Hasta ahora hemos trabajado con números naturales y con la adición. Dentro de \(\mathbb{N}\), sumar dos números no presenta ningún inconveniente: la suma de dos naturales siempre vuelve a ser un número natural.

Sin embargo, aparece una dificultad cuando intentamos realizar ciertas operaciones que habitualmente llamamos restas. Por ejemplo,
\[
7-3=4
\]
no parece presentar ningún problema, pues \(4\in\mathbb{N}\). Pero si intercambiamos los números,
\[
3-7,
\]
ya no encontramos ningún número natural que pueda representar esa cantidad.

El conjunto \(\mathbb{N}\) resulta demasiado pequeño para expresar todas las cantidades que queremos considerar. Necesitamos ampliarlo.

\subsection{Los números opuestos}

Consideremos el número \(3\). Queremos introducir un número que, al sumarse con \(3\), produzca \(0\). A este número lo llamamos \emph{opuesto} de \(3\) y lo escribimos
\[
-3.
\]
Por definición de opuesto,
\[
3+(-3)=0.
\]
De manera general, el opuesto de un número \(a\) se representa mediante
\[
-a
\]
y cumple
\[
a+(-a)=0.
\]
El signo \(-\) que aparece en \(-a\) no representa aquí una resta entre dos números. Indica que estamos tomando el \emph{opuesto} de \(a\).

Por ejemplo,
\[
-5,\qquad -18,\qquad -302
\]
son números, de la misma manera que
\[
5,\qquad18,\qquad302
\]
también lo son.
Observe además que el opuesto de \(0\) es el propio \(0\):
\[
-0=0,
\]
pues
\[
0+0=0.
\]

\subsection{Los números enteros}

Los números negativos que acabamos de introducir no pertenecen al conjunto de los números naturales. Por ejemplo,
\[
-3\notin\mathbb{N}.
\]
donde el símbolo \(\notin\) significa que no pertenece.

Para trabajar conjuntamente con los números naturales, sus opuestos y el cero utilizamos el conjunto de los números enteros, representado mediante \(\mathbb{Z}\):
\[
\mathbb{Z}
=
\{\ldots,-3,-2,-1,0,1,2,3,\ldots\}.
\]
De este modo, los números naturales son un subconjunto de los enteros
\[
\mathbb{N}\subseteq\mathbb{Z}.
\]
Es decir, los números
\[
1,2,3,\ldots
\]
son los enteros positivos (y también naturales); los números
\[
-1,-2,-3,\ldots
\]
son los enteros negativos; y el \(0\) no es positivo ni negativo.

La adición puede ahora considerarse sobre este conjunto:
\[
+\colon\mathbb{Z}\times\mathbb{Z}\to\mathbb{Z}.
\]
Esto significa que podemos adicionar dos números enteros cualesquiera y obtener nuevamente un número entero.

Por ejemplo,
\[
5+(-2),\qquad
(-4)+7,\qquad
(-8)+(-3)
\]
son todas expresiones perfectamente válidas de adición entre números enteros.

\subsection{Dos usos del símbolo menos}

Llegados a este punto conviene hacer una distinción importante. El símbolo \(-\) aparece en matemáticas con dos funciones relacionadas, pero diferentes.

En la expresión
\[
-5,
\]
el símbolo \(-\) actúa sobre un único número: indica el opuesto de \(5\). En este caso podemos llamarlo \emph{signo menos} u \emph{operador de oposición}.

En cambio, en
\[
8-5,
\]
el símbolo aparece entre dos números e indica una sustracción.

Ambos usos están íntimamente relacionados. De hecho, la sustracción puede definirse utilizando únicamente la adición y el concepto de opuesto.

\subsection{La sustracción como suma del opuesto}

Dados dos números enteros \(a\) y \(b\), definimos
\[
a-b=a+(-b).
\]
Esta igualdad contiene la idea fundamental de la sustracción:
\[
\boxed{\text{Restar un número equivale a sumar su opuesto.}}
\]
Por ejemplo,
\[
8-5
\]
puede escribirse como
\[
8+(-5).
\]
Ambas expresiones representan exactamente el mismo número:
\[
8-5=8+(-5)=3.
\]
Del mismo modo,
\[
3-7=3+(-7)=-4.
\]
Ahora podemos comprender por qué la operación que antes no podía efectuarse dentro de los números naturales sí puede evaluarse dentro de los enteros:
\[
-4\in\mathbb{Z}.
\]
No hemos necesitado inventar un procedimiento completamente independiente de la adición. La sustracción queda definida mediante una operación que ya conocemos y el concepto de número opuesto.

\subsection{El signo acompaña al sumando}

Esta forma de escribir la sustracción resulta especialmente útil cuando aparecen varias cantidades.

Considere, por ejemplo,
\[
2-10.
\]
Utilizando la definición anterior podemos escribir
\[
2-10=2+(-10).
\]
Ahora vemos claramente cuáles son los dos sumandos:
\[
2
\qquad\text{y}\qquad
-10.
\]
El segundo sumando es el número \(-10\), no el número \(10\).

Como ahora tenemos una adición, podemos utilizar la propiedad conmutativa:
\[
2+(-10)=(-10)+2.
\]
Por lo tanto,
\[
2-10
=
2+(-10)
=
(-10)+2
=
-8.
\]
Aquí aparece una idea que será muy útil en cálculos posteriores: una vez que una expresión ha sido escrita como una suma de enteros, cada sumando conserva su propio signo.

Por ejemplo,
\[
7-4+2
\]
puede interpretarse como
\[
7+(-4)+2.
\]
Los sumandos son entonces
\[
7,\qquad -4,\qquad 2.
\]
Al trabajar de esta manera resulta más sencillo distinguir los números que participan de la operación y reorganizarlos cuando sea conveniente.

\begin{note}
No debe confundirse esto con afirmar que la sustracción es conmutativa.

Por ejemplo,
\[
7-4=3,
\]
mientras que
\[
4-7=-3.
\]
Lo que sí podemos reorganizar libremente son los sumandos una vez que la expresión ha sido escrita como una adición:
\[
7+(-4)=(-4)+7.
\]
\end{note}

\subsection{Restar un número negativo}

La definición
\[
a-b=a+(-b)
\]
también explica qué ocurre cuando el número que restamos ya es negativo.
Considere
\[
7-(-3).
\]
Restar \(-3\) significa sumar su opuesto:
\[
7-(-3)=7+\bigl(-(-3)\bigr).
\]
Pero el opuesto de \(-3\) es \(3\):
\[
-(-3)=3.
\]
Por lo tanto,
\[
7-(-3)=7+3=10.
\]
Esto explica la conocida regla según la cual ``dos signos menos se convierten en un más''. No se trata de una regla aislada que deba memorizarse: es una consecuencia de tomar el opuesto de un número que ya era negativo.

En general,
\[
-(-a)=a
\]
y, por lo tanto,
\[
a-(-b)=a+b.
\]

\subsection{Algunas consecuencias}

A partir de la adición y de los números opuestos podemos obtener inmediatamente varias relaciones útiles.

Para cualquier \(a\in\mathbb{Z}\),
\[
a+0=a,
\]
pues \(0\) es el elemento neutro de la adición.
Además,
\[
a+(-a)=0,
\]
pues \(-a\) es el opuesto de \(a\).

Como la sustracción se define mediante la suma del opuesto, tenemos
\begin{align*}
a-a=a+(-a)=0,\\
a-0=a+(-0)=a,
\end{align*}
y
\[
0-a=0+(-a)=-a.
\]
También,
\[
a-(-b)=a+b.
\]
Estas expresiones no necesitan aprenderse como reglas independientes. Todas ellas pueden deducirse a partir de una única idea:
\[
a-b=a+(-b).
\]

